\documentclass[11pt, a4paper]{article}

% ── Packages ──────────────────────────────────────────────────────────────────
\usepackage[margin=2.5cm]{geometry}
\usepackage{fontenc}
\usepackage{inputenc}
% microtype disabled: font expansion incompatibility
\usepackage{hyperref}
\usepackage{xcolor}
\usepackage{graphicx}
\usepackage{booktabs}
\usepackage{array}
\usepackage{enumitem}
\usepackage{titlesec}
\usepackage{fancyhdr}
\usepackage{amsmath}
\usepackage{amssymb}
\usepackage{listings}
\usepackage{caption}
\usepackage{multicol}
\usepackage{setspace}

% ── Colors ────────────────────────────────────────────────────────────────────
\definecolor{solana-purple}{HTML}{9945FF}
\definecolor{solana-green}{HTML}{14F195}
\definecolor{solana-cyan}{HTML}{22D3EE}
\definecolor{dark-bg}{HTML}{0F0F15}
\definecolor{dark-card}{HTML}{18181B}
\definecolor{text-muted}{HTML}{71717A}
\definecolor{text-secondary}{HTML}{A1A1AA}

% ── Hyperref setup ────────────────────────────────────────────────────────────
\hypersetup{
  colorlinks=true,
  linkcolor=solana-purple,
  urlcolor=solana-cyan,
  citecolor=solana-green,
  pdftitle={SolAA: Zero-Movement PQC-Ready Account Abstraction for Solana},
  pdfauthor={Jason Wang},
  pdfsubject={Post-Quantum Cryptography, Solana, Account Abstraction},
}

% ── Section formatting ────────────────────────────────────────────────────────
\titleformat{\section}
  {\large\bfseries\color{solana-purple}}
  {\thesection.}{0.75em}{}[\vspace{-0.5em}\rule{\linewidth}{0.4pt}\vspace{0.2em}]

\titleformat{\subsection}
  {\normalsize\bfseries\color{solana-cyan}}
  {\thesubsection}{0.5em}{}

\titleformat{\subsubsection}
  {\normalsize\itshape\color{text-secondary}}
  {\thesubsubsection}{0.5em}{}

% ── Header / footer ───────────────────────────────────────────────────────────
\pagestyle{fancy}
\fancyhf{}
\fancyhead[L]{\small\color{text-muted}SolAA Whitepaper}
\fancyhead[R]{\small\color{text-muted}v1.0 --- May 2026}
\fancyfoot[C]{\small\color{text-muted}\thepage}
\renewcommand{\headrulewidth}{0.2pt}
\renewcommand{\headrule}{\hbox to\headwidth{\color{text-muted}\leaders\hrule height \headrulewidth\hfill}}

% ── Callout box ───────────────────────────────────────────────────────────────
\usepackage{tcolorbox}
\tcbuselibrary{skins}

\newtcolorbox{calloutbox}{
  colback=solana-purple!5,
  colframe=solana-purple,
  boxrule=0.8pt,
  arc=4pt,
  left=10pt, right=10pt, top=8pt, bottom=8pt,
  before skip=12pt, after skip=8pt,
}

\newtcolorbox{greenbox}{
  colback=solana-green!5,
  colframe=solana-green,
  boxrule=0.8pt,
  arc=4pt,
  left=10pt, right=10pt, top=8pt, bottom=8pt,
  before skip=12pt, after skip=8pt,
}

% ── List style ────────────────────────────────────────────────────────────────
\setlist[itemize]{leftmargin=1.5em, itemsep=2pt, parsep=0pt}
\setlist[enumerate]{leftmargin=1.5em, itemsep=2pt, parsep=0pt}

% ── Document ─────────────────────────────────────────────────────────────────
\begin{document}

% ── Title page ────────────────────────────────────────────────────────────────
\begin{titlepage}
  \centering
  \vspace*{2cm}

  {\Huge\bfseries\color{solana-purple} Sol\color{solana-green}AA}

  \vspace{0.75cm}
  {\LARGE\bfseries Zero-Movement PQC-Ready\\Account Abstraction for Solana}

  \vspace{0.5cm}
  \rule{0.6\linewidth}{0.4pt}

  \vspace{1cm}
  {\large Jason Wang \quad \textit{Cryptography Researcher}}\\[0.3em]
  {\color{text-muted}\small \href{https://solaa.hfi.network}{solaa.hfi.network}}

  \vspace{0.75cm}
  {\color{text-muted} May 2026 \quad v1.0}

  \vspace{2cm}

  \begin{abstract}
    \noindent
    SolAA is a production-ready account abstraction framework for Solana that delivers
    three capabilities simultaneously: \textbf{(1)} zero-movement key migration to
    post-quantum cryptography (ML-DSA-44), \textbf{(2)} ZK-proof-authorized smart account
    execution without Ed25519 signatures, and \textbf{(3)} trustless cross-chain ownership
    proofs without bridges. Unlike every prior account abstraction proposal, SolAA requires
    zero asset movement, zero new wallet creation, and zero transaction cost for the security
    upgrade itself. The entire \$200B+ Solana ecosystem is addressable on day one. This paper
    describes the protocol design, on-chain architecture, cryptographic foundations, and
    business rationale.
  \end{abstract}

  \vfill

  \begin{calloutbox}
    \textbf{Live Demo:} \href{https://solaa.hfi.network/demo/}{solaa.hfi.network/demo/}\\
    \textbf{Research:} 7 arXiv preprints (arXiv:2511.20505 and related)\\
    \textbf{Deployed:} Solana Devnet --- Program \texttt{EgKrUBUsQjC7BZ7xJGNLkDPP5UnvQ1u9Ldx7uRThNmL5}
  \end{calloutbox}
\end{titlepage}

% ── TOC ───────────────────────────────────────────────────────────────────────
\tableofcontents
\newpage

% ─────────────────────────────────────────────────────────────────────────────
\section{Executive Summary}
% ─────────────────────────────────────────────────────────────────────────────

Crypto wallets today work like a house with one lock and no spare key. If someone steals
your key --- or if a future quantum computer can copy any key ever made --- your assets are
gone forever. There is no way to change the lock without moving to a new house.

SolAA changes this. It lets Solana users upgrade to quantum-resistant security \textit{without
moving a single token, creating a new wallet, or paying any migration fee.} The wallet
address stays the same. The on-chain history stays intact. But the underlying security becomes
mathematically unbreakable --- even by quantum computers.

Beyond individual users, SolAA introduces two capabilities that do not exist anywhere on
Solana today:

\begin{itemize}
  \item \textbf{Smart accounts} authorized by zero-knowledge proofs rather than private keys,
    enabling programmable, recoverable, quantum-safe execution.
  \item \textbf{Cross-chain identity proofs} that let a Solana wallet prove ownership of
    accounts on Ethereum, Tron, or any other blockchain --- without bridges, without fees,
    and without trusting any intermediary.
\end{itemize}

\begin{calloutbox}
  \textbf{The one-sentence pitch:} SolAA is the security upgrade layer that makes every
  existing Solana wallet quantum-safe and interoperable --- with no migration required.
\end{calloutbox}

% ─────────────────────────────────────────────────────────────────────────────
\section{The Problem}
% ─────────────────────────────────────────────────────────────────────────────

\subsection{Permanent Key Vulnerability}

Every Solana wallet today is controlled by a single Ed25519 private key. This creates
three compounding risks:

\begin{enumerate}
  \item \textbf{Irrecoverable compromise.} If a private key is leaked --- through malware,
    phishing, a compromised backup, or a supply chain attack --- the account is permanently
    lost. There is no mechanism to revoke the old key and issue a new one without moving
    all assets.

  \item \textbf{Quantum threat.} Shor's algorithm running on a sufficiently capable quantum
    computer can derive an Ed25519 private key from its public key in polynomial time.
    NIST finalized post-quantum cryptography standards in 2024 (FIPS 204 / ML-DSA).
    Financial regulators in the US, EU, and China are actively drafting mandates.

  \item \textbf{Single point of failure.} A seed phrase written on paper, stored in a
    password manager, or memorized is the only recovery path. Loss means permanent
    loss.
\end{enumerate}

\subsection{The Migration Tax}

Existing solutions to the above problems --- hardware wallets, multisig, EVM-style account
abstraction (ERC-4337) --- all share the same fundamental flaw: they require users to
\textit{migrate assets to a new address}. This creates:

\begin{itemize}
  \item Taxable events in most jurisdictions
  \item Loss of on-chain history (NFT provenance, DeFi reputation, governance records)
  \item High friction --- most users never do it
  \item Protocol-level complexity for dApps that track wallet addresses
\end{itemize}

\subsection{Cross-Chain Identity is Broken}

As of 2026, proving that a Solana wallet and an Ethereum wallet belong to the same person
requires either (a) bridging assets between chains --- costing gas, time, and carrying
\$2.5B+ in historical hack risk --- or (b) trusting a centralized attestation service.
Neither is acceptable for high-value use cases like institutional KYC, cross-chain DAO
voting, or multichain airdrop eligibility.

% ─────────────────────────────────────────────────────────────────────────────
\section{The SolAA Solution}
% ─────────────────────────────────────────────────────────────────────────────

SolAA solves all three problems through a single unified architecture with three layers.

\subsection{Layer 1: Zero-Movement Key Upgrade (SA-Migration)}

The user's existing seed phrase is encapsulated into a \textbf{REV32 Sealed Artifact}:
an AES-256-GCM-SIV encrypted blob derived via Argon2id (memory: 4 MB, iterations: 3)
from a user-supplied passphrase. The resulting \texttt{id\_com} (identity commitment) is
a 32-byte cryptographic commitment that:

\begin{itemize}
  \item Preserves the user's existing Solana address (no asset movement)
  \item Acts as the root of the post-quantum key hierarchy
  \item Is the seed for all downstream SolAA operations
\end{itemize}

\begin{greenbox}
  \textbf{Key property:} The identity commitment is derived from the user's existing
  mnemonic, so the \textit{same wallet address} is preserved. Zero transactions are
  broadcast. Zero gas is spent. The upgrade is entirely client-side until the user
  explicitly creates an on-chain smart account.
\end{greenbox}

\subsection{Layer 2: ZK-Authorized Smart Account}

A Program Derived Address (PDA) is created on-chain with the following properties:

\begin{center}
\begin{tabular}{ll}
  \toprule
  \textbf{Property} & \textbf{Value} \\
  \midrule
  Address derivation & \texttt{["solaa", id\_com[0..32]]} \\
  Authorization & Groth16 ZK proof (128 bytes) \\
  Verification cost & $\approx$170,000 compute units \\
  Key rotation & Supported, same address retained \\
  Recovery delay & 1,512,000 slots ($\approx$7 days) \\
  Nonce & Replay-proof, increments per instruction \\
  \bottomrule
\end{tabular}
\end{center}

\vspace{0.5em}
The account is authorized by a \textbf{ZK-ACE proof} (zero-knowledge authorization circuit)
rather than an Ed25519 signature. This means:

\begin{itemize}
  \item The private key never touches the transaction
  \item Authorization is programmable (time-locks, spending limits, social recovery)
  \item A 128-byte Groth16 proof replaces a 2,420-byte ML-DSA-44 signature on the wire
\end{itemize}

\subsubsection{Key Rotation}

Users can rotate to a new ML-DSA-44 (post-quantum) key at any time. The rotation instruction:

\begin{enumerate}
  \item Verifies a ZK proof of knowledge of the old \texttt{id\_com}
  \item Replaces the on-chain authority with the new PQC public key
  \item Bumps the account nonce to invalidate any pending operations
  \item Preserves the PDA address --- no migration required
\end{enumerate}

\subsection{Layer 3: Cross-Chain Ownership Proof (ZK-Ownership)}

SolAA generates a \textbf{ZK-STARK proof} that encodes:

\begin{itemize}
  \item The user's Solana \texttt{id\_com}
  \item A target foreign address (e.g., an Ethereum \texttt{0x...} address)
  \item A foreign chain ID
\end{itemize}

All three are \textit{public inputs} to the circuit. Any verifier --- on any chain that
supports the ZK-ACE verifier contract --- can confirm the proof without learning the
user's private key or seed phrase.

\begin{calloutbox}
  \textbf{What this enables:} Cross-chain airdrops, multichain DAO voting, unified KYC,
  and interoperable DeFi reputation --- all from a single Solana wallet, with no bridges,
  no wrapped assets, and no trusted intermediaries.
\end{calloutbox}

% ─────────────────────────────────────────────────────────────────────────────
\section{Technical Architecture}
% ─────────────────────────────────────────────────────────────────────────────

\subsection{On-Chain Program}

The SolAA Anchor program (deployed at \texttt{EgKrUBUs\ldots}) exposes five instructions:

\begin{center}
\begin{tabular}{lp{9cm}}
  \toprule
  \textbf{Instruction} & \textbf{Description} \\
  \midrule
  \texttt{create\_account} & Initializes a PDA smart account from an \texttt{id\_com} \\
  \texttt{execute} & Executes an authorized CPI payload given a valid ZK proof \\
  \texttt{rotate\_key} & Rotates the PQC authority key, preserving the PDA address \\
  \texttt{initiate\_recovery} & Begins time-locked VA-DAR recovery process \\
  \texttt{verify\_ownership} & Records a cross-chain ownership proof on-chain \\
  \bottomrule
\end{tabular}
\end{center}

\subsection{ZK-ACE Circuit}

The authorization circuit is implemented in Stwo (Rust), a STARK-based proving system.
For the demo, proofs are generated server-side via the Axum API and verified on-chain
via the alt\_bn128 pairing precompile. Circuit statistics:

\begin{itemize}
  \item Constraints: $\approx$4,024 R1CS
  \item Proof size: 128 bytes (Groth16)
  \item Proving time: $\approx$52 ms (server-side)
  \item On-chain verification: $\approx$170,000 CU
\end{itemize}

\subsection{Cryptographic Primitives}

\begin{center}
\begin{tabular}{lll}
  \toprule
  \textbf{Component} & \textbf{Primitive} & \textbf{Standard} \\
  \midrule
  Key derivation & Argon2id & RFC 9106 \\
  Symmetric encryption & AES-256-GCM-SIV & RFC 8452 \\
  PQC signature & ML-DSA-44 & NIST FIPS 204 \\
  ZK proof system & Groth16 / STARK & Stwo \\
  Hash & SHA-256 & FIPS 180-4 \\
  Elliptic curve pairing & alt\_bn128 & Solana precompile \\
  \bottomrule
\end{tabular}
\end{center}

% ─────────────────────────────────────────────────────────────────────────────
\section{Business Case}
% ─────────────────────────────────────────────────────────────────────────────

\subsection{Market Size}

\begin{itemize}
  \item \textbf{\$200B+ in Solana assets} can upgrade to quantum-safe security without
    any migration cost. The addressable market is every existing Solana user.
  \item \textbf{600M+ crypto users globally} will eventually need PQC-safe wallets as
    regulations tighten. SolAA is the only zero-friction path on Solana.
  \item \textbf{Cross-chain identity} is a \$50B+ opportunity currently served only by
    fragile bridge infrastructure.
\end{itemize}

\subsection{Why Now}

\begin{enumerate}
  \item \textbf{NIST finalized PQC standards (2024).} The regulatory clock is running.
    Financial institutions are already receiving guidance to begin PQC migration.
  \item \textbf{Bridge hacks continue.} \$2.5B+ lost since 2021. Every high-profile
    bridge hack increases demand for bridge-free cross-chain solutions.
  \item \textbf{Account abstraction is a hot category.} ERC-4337 proved the demand.
    SolAA brings it to Solana with a materially better migration story.
\end{enumerate}

\subsection{Competitive Differentiation}

\begin{center}
\begin{tabular}{lccc}
  \toprule
  \textbf{Feature} & \textbf{SolAA} & \textbf{ERC-4337} & \textbf{Traditional Multisig} \\
  \midrule
  Zero asset migration & \checkmark & $\times$ & $\times$ \\
  PQC-ready (day one) & \checkmark & $\times$ & $\times$ \\
  Key rotation (same addr) & \checkmark & $\times$ & $\times$ \\
  Cross-chain identity & \checkmark & $\times$ & $\times$ \\
  No bridge required & \checkmark & N/A & N/A \\
  Solana-native & \checkmark & $\times$ & Partial \\
  \bottomrule
\end{tabular}
\end{center}

% ─────────────────────────────────────────────────────────────────────────────
\section{Research Foundation}
% ─────────────────────────────────────────────────────────────────────────────

Every design decision in SolAA is grounded in published cryptographic research.
The following papers form the theoretical basis of the protocol:

\begin{center}
\begin{tabular}{lll}
  \toprule
  \textbf{Paper} & \textbf{Reference} & \textbf{Contribution} \\
  \midrule
  ACE-GF & arXiv:2511.20505 & Identity derivation framework \\
  ZK-ACE & arXiv:2603.07974 & ZK authorization for PQC \\
  AR-ACE & arXiv:2603.07982 & Proof-off-path relay architecture \\
  VA-DAR & arXiv:2603.02690 & Decentralized address recovery \\
  CT-DAP & arXiv:2603.07933 & Destroyable authorization paths \\
  ACE Runtime & arXiv:2603.10242 & Sub-second finality blockchain runtime \\
  \bottomrule
\end{tabular}
\end{center}

% ─────────────────────────────────────────────────────────────────────────────
\section{Roadmap}
% ─────────────────────────────────────────────────────────────────────────────

\begin{enumerate}
  \item \textbf{Q2 2026 --- Devnet launch (current).}
    Full demo live at \href{https://solaa.hfi.network/demo/}{solaa.hfi.network/demo/}.
    All five on-chain instructions deployed and functional.

  \item \textbf{Q3 2026 --- SDK \& integrations.}
    Open-source TypeScript SDK. Integration guides for major Solana dApps.
    WASM-compiled ZK prover for client-side proof generation.

  \item \textbf{Q4 2026 --- Mainnet beta.}
    Security audit. Mainnet deployment. First institutional partners.

  \item \textbf{2027 --- Cross-chain verifier rollout.}
    Deploy ZK-ACE verifier contracts on Ethereum and Tron.
    Enable live cross-chain identity proofs.

  \item \textbf{2027+ --- Protocol governance.}
    Decentralized upgrade authority. Community governance for circuit upgrades.
\end{enumerate}

% ─────────────────────────────────────────────────────────────────────────────
\section{Conclusion}
% ─────────────────────────────────────────────────────────────────────────────

SolAA is not a theoretical proposal. It is a working system, live on Solana Devnet,
with a full-stack demo, an Anchor on-chain program, a Rust API server, and a research
foundation built on seven published cryptographic papers.

The window for establishing the PQC-safe account abstraction standard on Solana is
open now --- before regulatory mandates force a chaotic, expensive migration. SolAA
offers the only path that upgrades every existing user with zero friction.

\begin{greenbox}
  \textbf{Try the demo:} \href{https://solaa.hfi.network/demo/}{solaa.hfi.network/demo/}\\
  \textbf{Contact:} \href{mailto:zwanjas@gmail.com}{zwanjas@gmail.com}
\end{greenbox}

\vspace{2em}
\noindent\color{text-muted}\small
This document is provided for informational purposes. Protocol parameters and timelines
are subject to change. Nothing herein constitutes financial or legal advice.

\end{document}
